\documentclass[a4,11pt]{aleph-notas}
% Se puede ver la documentación aquí:
% https://github.com/alephsub0/LaTeX_aleph-notas

% -- Paquetes adicionales
\usepackage{enumitem}
\usepackage{url}
\usepackage{array}
\usepackage{booktabs}
\usepackage{longtable}
\hypersetup{
    urlcolor=blue,
    linkcolor=blue,
}

% -- Datos
\institucion{Facultad de Ciencias Exactas, Naturales y Ambientales}
\carrera{Catálogo STEM}
\asignatura{Álgebra lineal}
\tema{Clase 09: Espacios con producto interno}
\autor{Andrés Merino}
\fecha{Periodo 2026-01}

\logouno[0.16\textwidth]{Logos/logoPUCE_04_ac}
\definecolor{colortext}{HTML}{1957d1}
\definecolor{colordef}{HTML}{1957d1}
\fuente{montserrat}

\begin{document}

\encabezado 

%%%%%%%%%%%%%%%%%%%%%%%%%%%%%%%%%%%%%%%%
\section*{Resultado de Aprendizaje}
%%%%%%%%%%%%%%%%%%%%%%%%%%%%%%%%%%%%%%%%

%%%%%%%%%%%%%%%%%%%%%%%%%%%%%%%%%%%%%%%%
\subsection*{RdA de la asignatura:}
\begin{itemize}[leftmargin=*]
    \item \textbf{RdA 1:} Comprender los conceptos fundamentales del Álgebra Lineal, incluyendo el estudio de matrices, determinantes, sistemas de ecuaciones lineales y espacios vectoriales, destacando su importancia en el análisis y resolución de problemas matemáticos.
    \item \textbf{RdA 2:} Resolver problemas prácticos y teóricos relacionados con el Álgebra Lineal, utilizando técnicas de cálculo con matrices, determinantes y sistemas de ecuaciones, así como en el análisis de espacios vectoriales.    
    % \item \textbf{RdA 3:} Aplicar los principios y métodos del Álgebra Lineal en una variedad de contextos demostrando comprensión de su relevancia y utilidad en situaciones prácticas y teóricas.
\end{itemize}


%%%%%%%%%%%%%%%%%%%%%%%%%%%%%%%%%%%%%%%%
\subsection*{RdA de la actividad:}
\begin{itemize}[leftmargin=*]
    \item Calcular productos internos, normas y distancias en distintos espacios vectoriales.
    \item Determinar ortogonalidad de vectores en $\mathbb{R}^n$, polinomios y funciones.
    \item Aplicar Gram-Schmidt para construir bases ortogonales y ortonormales.
    \item Determinar complementos ortogonales y calcular proyecciones ortogonales.
\end{itemize}

%%%%%%%%%%%%%%%%%%%%%%%%%%%%%%%%%%%%%%%%
\section*{Introducción}
%%%%%%%%%%%%%%%%%%%%%%%%%%%%%%%%%%%%%%%%

%%%%%%%%%%%%%%%%%%%%%%%%%%%%%%%%%%%%%%%%
\paragraph{Control de lectura:}
Antes de iniciar la clase, se aplicará el control de lectura del siguiente enlace:
\href{https://gemini.google.com/share/7d22976ac15d}{Control de lectura - Clase 09}.

\paragraph{Pregunta inicial:}
¿Dos polinomios pueden ser perpendiculares? ¿Y dos matrices?

%%%%%%%%%%%%%%%%%%%%%%%%%%%%%%%%%%%%%%%%
\section*{Desarrollo}
%%%%%%%%%%%%%%%%%%%%%%%%%%%%%%%%%%%%%%%%

%%%%%%%%%%%%%%%%%%%%%%%%%%%%%%%%%%%%%%%%
\subsection*{Actividad 1: Producto interno, norma y ortogonalidad}

La clase inicia con una recuperación breve de norma y distancia en $\mathbb{R}^n$. Luego, mediante clase magistral y práctica guiada, se introduce la noción de producto interno en espacios abstractos, sus propiedades, la ortogonalidad y resultados asociados como Pitágoras, Cauchy-Schwarz y desigualdad triangular.

\paragraph{¿Cómo lo haremos?}
\begin{itemize}[leftmargin=*]
    \item \textbf{Clase magistral:} se presentan las definiciones de producto interno, norma, distancia, ortogonalidad, conjuntos ortogonales/ortonormales, complemento ortogonal y proyección ortogonal. Se usan ejemplos en $\mathbb{R}^n$, matrices, polinomios y funciones. Se usa el resumen \href{https://fcena-puce.github.io/AlgebraLineal-05-N0048/01Resumen/Resumen09.pdf}{Resumen09.pdf}.
    \item \textbf{Resolución de ejercicios:} se guiará a los estudiantes en la solución de ejercicios del documento \href{https://fcena-puce.github.io/AlgebraLineal-05-N0048/04Ejercicios/Ejercicios09.pdf}{Ejercicios09.pdf}, calculando productos internos, verificando ortogonalidad, aplicando Gram-Schmidt y calculando proyecciones.
    \item \textbf{Recomendaciones de lectura:} 
    \begin{itemize}
        \item \href{https://blog.nekomath.com/algebra-lineal-i-producto-interior-y-desigualdad-de-cauchy-schwarz/}{Producto interior y desigualdad de Cauchy-Schwarz}
    \end{itemize}
\end{itemize}

\paragraph{Verificación de aprendizaje:}
\begin{enumerate}[leftmargin=*]
    \item ¿Qué propiedades debe cumplir una aplicación para ser un producto interno?
    \item ¿Cómo se determina si dos vectores son ortogonales en un espacio de polinomios o funciones?
    \item ¿Cómo se usa Gram-Schmidt para construir una base ortonormal?
    \item ¿Qué interpretación geométrica tiene la proyección ortogonal?
\end{enumerate}

%%%%%%%%%%%%%%%%%%%%%%%%%%%%%%%%%%%%%%%%
\section*{Cierre}
%%%%%%%%%%%%%%%%%%%%%%%%%%%%%%%%%%%%%%%%

\paragraph{Tarea:}
Resolver del libro \href{https://catalogobiblioteca.puce.edu.ec/cgi-bin/koha/opac-detail.pl?biblionumber=86081}{Fundamentos de álgebra lineal de Larson}: sección 5.2, ejercicios: 1,3,5,7,9,17,21,31,32.39

\paragraph{Pregunta de investigación:}
\begin{enumerate}[leftmargin=*]
    \item ¿Qué ventajas tiene trabajar con bases ortonormales frente a bases arbitrarias?
    \item ¿Cómo se relaciona el complemento ortogonal con la solución de problemas de aproximación?
    \item ¿Por qué la proyección ortogonal entrega el vector más cercano en un subespacio?
\end{enumerate}

\paragraph{Para la próxima clase:}  
Visualizar el video \href{https://youtu.be/YJfS4_m_0Z8?si=kFS74G3kXhWGWq4L}{Transformaciones lineales}.
Realizar la actividad de clase invertida: \href{https://fcena-puce.github.io/AlgebraLineal-05-N0048/07Act-ClaseInvertida/03ClsInv-AplicacionLineal.pdf}{03ClsInv-AplicacionLineal.pdf}.



\end{document}
